\section{Shared preparation, invariant polytopes, and noncyclic memory}
\label{sec:r77-shared-preparation}

A device that reproduces one prepared response need not reproduce a family
of preparations with the same memory, transitions, and decoder.  We now
solve the structural problem with this common-device requirement.  The
preparation is selected before the command word; it is not revealed to the
terminal decoder.  In particular, one may not choose a different decoder
after learning which preparation was selected.  This convention is what
makes joint realization a resource question.

Fix a finite group $G$, a surjective command homomorphism
$\gamma:A^*\to G$, a finite outcome set $C$, and a nonempty finite
preparation set $I$.  A \emph{shared controller} consists of one family of
stochastic matrices $Q_a$ and one stochastic decoder $D_0$, together with
an initial row for each $i\in I$.  In the pure convention these rows are
prescribed label rows $e_{o_i}$; different preparations may use the same
label when their required laws agree.  All labels, including initialization
labels, are counted.  In the mixed convention the initial rows may be
arbitrary probability rows and their external preparation is not counted.
The result concerns this explicit distinction, not the bit cost of selecting
$i$ or an uncounted label made available to the decoder.

For each $i,g$ let $K_{i,g}\subseteq\Delta(C)$ be closed and convex,
possibly empty.  The requirement is
\begin{equation}\label{eq:r77-shared-specification}
 e_{o_i}Q_wD_0\in K_{i,\gamma(w)}
 \qquad(i\in I,\ w\in A^*),
\end{equation}
with arbitrary initial rows in the mixed convention.  For probability
rows $\Lambda=(\lambda_i)_{i\in I}$ on a finite set put
\[
 \iota(\Lambda)=
 \bigl|\{\lambda_i:i\in I,\ \lambda_i\hbox{ is not a point mass}\}\bigr|.
\]
Repeated rows are counted once.

\begin{theorem}[Shared normal form and exact initialization charge]
\label{thm:r77-shared-normal-form}
The specification \eqref{eq:r77-shared-specification} has a pure shared
controller on at most $M$ labels if and only if there are a finite
$G$-set $X$, probability rows $\lambda_i$ on $X$, and one stochastic
decoder $D$ such that
\begin{equation}\label{eq:r77-shared-core}
 \lambda_iP_gD\in K_{i,g}\quad(i\in I,g\in G),
 \qquad |X|+\iota(\Lambda)\le M.
\end{equation}
For mixed preparations the criterion is $|X|\le M$.  The action may have
a kernel.  Thus a single permutation core suffices for all preparations,
and each distinct mixed preparation costs exactly one additional label
in this normal form.
\end{theorem}
\begin{proof}
Use the identity-word construction in the proof of
Theorem~\ref{thm:r76-convex-normal-form}.  Its averaging matrix $T$,
stochastic idempotent $E=\lim_nT^n$, and disjoint closed classes
$C_1,\ldots,C_k$ depend only on the command matrices, not on the initial
row.  Write $\mu_x$ for their stationary probability rows.  The sandwiched
maps $EQ_uE$ permute these rows and give one action $P_{\gamma(u)}$ on
$X=\{1,\ldots,k\}$.  Put
\[
 e_{o_i}E=\sum_{x\in X}\lambda_i(x)\mu_x,
 \qquad D(x,\cdot)=\mu_xD_0.
\]
The same limit as in that theorem gives
\[
 \lambda_iP_gD
   =\lim_{n\to\infty}e_{o_i}T^nQ_uT^nD_0\in K_{i,g},
 \qquad \gamma(u)=g,
\]
because each approximant is a convex combination of laws for words
representing $g$.  There is no averaging between different preparations.

If $o_i$ lies in a closed class, then $e_{o_i}E$ is its stationary row,
so $\lambda_i$ is a point mass.  Distinct mixed $\lambda_i$ must therefore
arise from distinct initial labels outside the closed classes.  These
classes use at least $k$ labels, so
$\iota(\Lambda)\le M-k$.  This argument also permits an outside initial
label to project to a point mass, and it does not charge that label when
the projected realization no longer needs it.  It permits several initial
labels to have the same projected mixed row, in which case they are merged.
For mixed initial rows the same projection gives $k\le M$ without the
additional count.

Conversely, retain the permutation labels $X$.  For every distinct
nonpoint row $\lambda$ among the $\lambda_i$, adjoin one label $\star_\lambda$
with decoder $\lambda D$ and, on command $a$, transition row
$\lambda P_{\gamma(a)}$ into $X$.  No transition returns to an adjoined
label.  A point row uses its own point of $X$ as initial label; a nonpoint
row uses $\star_\lambda$.  The empty-word law is $\lambda_iD$ and the law
after every nonempty word is $\lambda_iP_{\gamma(w)}D$.  This establishes
\eqref{eq:r77-shared-core} with the exact displayed count.  The mixed
case starts directly from $\lambda_i$ on $X$.
\end{proof}

For exact targets write $\mathcal B=\{b_i:i\in I\}$, where
$b_i:G\to\Delta(C)$; duplicate target tables are identified.
Let $\kappa_G(\mathcal B)$ and $m_G(\mathcal B)$ denote the mixed and pure
shared minima.  The first equals the least $|X|$ in a simultaneous
factorization $b_i(g)=\lambda_iP_gD$; the second includes
$\iota(\Lambda)$.  In particular
\[
 \max_i\kappa_G(b_i)\le\kappa_G(\mathcal B),
 \qquad
 \kappa_G(\mathcal B)\le m_G(\mathcal B).
\]
Neither inequality compares a collection of separately optimized decoders
with a single common decoder.

\subsection{A geometric formula with all initialization labels counted}
Let $\mathcal P_G(\mathcal B)$ be the set of nonempty polytopes in the
finite-dimensional product simplex
\[
 \mathcal Q=\{q:G\to\Delta(C)\}
\]
that contain $\mathcal B$ and are invariant under
$(q\cdot h)(g)=q(hg)$.  The whole product simplex is one such polytope.
Write $v(P)=|\operatorname{vert}P|$.  The polytope is allowed to extend
beyond the linear span of the prescribed responses.  Forbidding this
extension would impose an additional positive-realization constraint.

\begin{theorem}[Invariant-polytope resource formulas]
\label{thm:r77-shared-polytope}
For every nonempty finite target family,
\begin{align}
 \kappa_G(\mathcal B)
   &=\min_{P\in\mathcal P_G(\mathcal B)}v(P),
       \label{eq:r77-polytope-mixed}\\
 m_G(\mathcal B)
   &=\min_{P\in\mathcal P_G(\mathcal B)}
       \left(v(P)+|\mathcal B\setminus\operatorname{vert}P|\right).
       \label{eq:r77-polytope-pure}
\end{align}
The minima are attained.  Every minimizing polytope constructs an optimal
controller: its vertices are permutation labels, and each distinct
prescribed nonvertex response receives one transient initialization label.
\end{theorem}
\begin{proof}
Given a simultaneous factorization, associate to each core label the
profile $q_x(g)=e_xP_gD$.  The set of profiles is invariant under $G$,
and $P=\operatorname{conv}\{q_x:x\in X\}$ belongs to
$\mathcal P_G(\mathcal B)$.  Every vertex is a profile, so $v(P)\le|X|$.
This proves the necessary mixed bound.

For the pure bound, apply the preceding normal form to an arbitrary
$M$-label pure realization.  Let $s=\iota(\Lambda)$ and $k=|X|$, so
$k+s\le M$.  If a target $b_i$ is not a vertex of $P$, it either equals
a nonvertex core profile $q_x$, or is not any core profile.  There are
at most $k-v(P)$ distinct profiles of the first kind: at least $v(P)$
core labels are required to represent the different vertices.  Every
distinct target of the second kind requires a distinct nonpoint
$\lambda_i$, so there are at most $s$ such targets.  Therefore
\[
 v(P)+|\mathcal B\setminus\operatorname{vert}P|
       \le v(P)+(k-v(P))+s\le M.
\]
This is the step that prevents a nonextreme continuing label used by
another preparation from being charged twice.

Conversely, a finite group preserving $P$ permutes its vertices.  Use
these vertices as $X$, with $D(q,c)=q(e,c)$.  Barycentric coordinates
of $b_i\in P$ supply a probability row $\lambda_i$ with
$\lambda_iP_gD=b_i(g)$.  If $b_i$ is a vertex choose its point row.
If it is not a vertex, every barycentric representation is nonpoint.
Choose the same representation for identical target tables.  Different
target tables have different representations, since their barycentric
images differ.  The normal-form construction thus costs exactly
$v(P)+|\mathcal B\setminus\operatorname{vert}P|$ in the pure convention,
and $v(P)$ in the mixed convention.

The feasible integer costs form nonempty subsets of the positive
integers.  Their least elements are realized by some feasible polytope;
the two constructions prove the asserted equalities and attainment.
\end{proof}

Invariant cones in positive realization and the relation between symmetry
and nonnegative factorizations are classical
\cite{Heller65,BenvenutiFarina04,GPT13}.  Formula
\eqref{eq:r77-polytope-mixed} is a response-polytope form of that structure.
The additional resource statement here is the exact shared-device
pure-initialization term in \eqref{eq:r77-polytope-pure}, obtained after
the common identity-word projection.  It neither assumes that a supplied
stochastic controller satisfies the group relations nor identifies linear
rank with positive memory.

\subsection{Joint Fourier multiplicity and exact low-contrast synthesis}
Let
\[
 V_{\mathcal B}=\operatorname{span}_{\mathbb R}
    \{b_i\cdot h:i\in I,h\in G\},\qquad d=\dim V_{\mathcal B}.
\]
For an irreducible unitary representation $\rho$ of degree $d_\rho$,
use the stacked matrix $A_\rho(b_i)$ from
\eqref{eq:r75-fourier-block} and define the \emph{joint} block
\begin{equation}\label{eq:r77-joint-fourier}
 \mathcal A_\rho(\mathcal B)=
   \bigl[A_\rho(b_1)\ \cdots\ A_\rho(b_{|I|})\bigr],
 \qquad r_\rho(\mathcal B)=\operatorname{rank}_{\mathbb C}
                                  \mathcal A_\rho(\mathcal B).
\end{equation}
Unlike the single-preparation rank, $r_\rho(\mathcal B)$ can exceed
$d_\rho$.  The same decoder must accommodate the union of the Fourier
column spaces, rather than their largest separate rank.

\begin{lemma}[Joint multiplicity]
\label{lem:r77-joint-multiplicity}
The multiplicity of $\rho$ in
$V_{\mathcal B}\otimes_{\mathbb R}\mathbb C$ equals
$r_\rho(\mathcal B)$.  Consequently every simultaneous factorization
on a $G$-set $X$ satisfies
\[
 r_\rho(\mathcal B)\le m_\rho(X)\qquad(\rho\in\widehat G).
\]
\end{lemma}
\begin{proof}
Fourier transformation sends the translate $b_i\cdot h$ to
$A_\rho(b_i)\rho(h)$.  The group-algebra decomposition used in
Lemma~\ref{lem:r75-fourier-multiplicity} therefore identifies its
$\rho$-component with
\[
 \left\{\sum_{i\in I}A_\rho(b_i)M_i:
                       M_i\in M_{d_\rho}(\mathbb C)\right\}.
\]
The matrices $M_i$ are independent, since the different generators may
have independent group-algebra coefficients.  This is the space of all
matrices whose columns belong to the sum of the column spaces of the
$A_\rho(b_i)$.  That sum has dimension $r_\rho(\mathcal B)$, so the
right module consists of precisely that many rows of length $d_\rho$.
For a simultaneous positive factorization, the span of all reachable
preparation rows is a submodule of $\mathbb R^X$, and evaluation against
the common decoder maps it onto $V_{\mathcal B}$.  Complete reducibility
therefore embeds $V_{\mathcal B}$ into $\mathbb R^X$, proving the bound.
\end{proof}

Write $\nu_G(V_{\mathcal B})$ for the smallest size of a permutation
module containing $V_{\mathcal B}$.  For one representative $H$ of each
subgroup conjugacy class, Frobenius reciprocity gives
\begin{equation}\label{eq:r77-joint-module-cover}
 \nu_G(V_{\mathcal B})=
 \min\left\{\sum_Hq_H[G:H]:q_H\in\mathbb Z_{\ge0},\quad
       \sum_Hq_Hm_\rho(H)\ge r_\rho(\mathcal B)
                           \ (\rho\in\widehat G)\right\}.
\end{equation}
The equality follows by orbit decomposition and the real-descent part
of Lemma~\ref{lem:r75-fourier-multiplicity}; it applies also to complex
and quaternionic representation types.

\begin{theorem}[Exact noncyclic memory at small common contrast]
\label{thm:r77-shared-low-contrast}
Suppose that all preparations have the same group-averaged law,
\[
 \frac1{|G|}\sum_{g\in G}b_i(g)=\pi\qquad(i\in I),
\]
and put $b_{i,t}(g)=\pi+t\{b_i(g)-\pi\}$ for $0\le t\le1$.
There is an explicitly positive $t_0$ such that
\begin{equation}\label{eq:r77-shared-low-contrast}
 \kappa_G(\{b_{i,t}:i\in I\})
      =\nu_G(V_{\mathcal B})\qquad(0<t\le t_0).
\end{equation}
For every $t>0$ the joint module and its Fourier multiplicities are
unchanged.  The exact pure resource is the invariant-polytope formula
\eqref{eq:r77-polytope-pure}; in particular, if $r=|\mathcal B|$ and
$N=|\{b_i\cdot g:i\in I,g\in G\}|$, then on this interval
\[
 \nu_G(V_{\mathcal B})\le m_G(\{b_{i,t}\})
       \le\min\{N,\nu_G(V_{\mathcal B})+r\}.
\]
At $t=0$ both resources equal one.
\end{theorem}
\begin{proof}
Remove outcomes with $\pi_c=0$; positivity makes them identically zero
for every preparation and group element.  They will be restored as
zero decoder columns.  Let $\overline b$ be the constant table $\pi$.
Group averaging maps every generating residual to $\overline b$, hence
\[
 V_{\mathcal B}=\mathbb R\overline b\oplus W,
 \qquad W=\ker\left(|G|^{-1}\sum_{g\in G}L_g\right).
\]
The nontrivial part $W$ has no invariant vector and all its elements have
zero total output mass.  If $W=0$, all targets equal $\overline b$ and
the assertion is immediate.  Otherwise choose a minimum $G$-set $X$
with $k=\nu_G(V_{\mathcal B})$ and a real equivariant injection
$j:W\to\mathbb R^X$.  Let $F:\mathbb R^X\to W$ be orthogonal
projection onto $j(W)$ followed by $j^{-1}$.  Since the permutation
inner product is invariant, $F$ is equivariant and annihilates every
invariant row.  Write
\[
 p=(1/k,\ldots,1/k),\qquad
 a_i=j(b_i-\overline b),\qquad
 T(x,c)=F(e_x)(e,c).
\]
Then $a_i\mathbf1=0$, $T\mathbf1=0$, $pT=0$, and
$a_iP_gT=b_i(g)-\pi$.  Put
\[
 A=\max_i\|a_i\|_\infty>0,\quad
 \epsilon=\frac1{2kA},\quad
 B=\max_{x,c}\frac{|T(x,c)|}{\pi_c}>0,\quad
 t_0=\min\{1,\epsilon/B\}.
\]
The rows $\lambda_i=p+\epsilon a_i$ have all entries at least
$1/(2k)$.  For $0\le t\le t_0$ the \emph{common} matrix
$D_t(x,c)=\pi_c+(t/\epsilon)T(x,c)$ is stochastic, and
\[
 \lambda_iP_gD_t=\pi+t\{b_i(g)-\pi\}=b_{i,t}(g)
\]
for every preparation simultaneously.  This proves the matching upper
bound, with one decoder and one action.

The equivariant map that is the identity on $\mathbb R\overline b$
and multiplication by $t$ on $W$ is invertible for $t>0$ and sends all
residuals to their contrasted versions.  Thus their joint module and
all equality relations are unchanged.  Lemma~\ref{lem:r77-joint-multiplicity}
and \eqref{eq:r77-joint-module-cover} give the lower bound.  The constructed
preparations have at most $r$ distinct nonpoint rows, so their pure
initialization costs at most $k+r$.  Alternatively, the union of all
residuals is a deterministic shared machine with $N$ labels.  At zero
contrast one common constant decoder and one fixed label suffice.
\end{proof}

The common-mean hypothesis has mathematical content.  Distinct constant
preparation laws can already have positive-factorization cost larger
than their linear rank; contrast does not change them.  For example,
with the trivial group, the four preparation rows
\[
 \frac12\begin{pmatrix}
 1&1&0&0\\0&1&1&0\\0&0&1&1\\1&0&0&1
 \end{pmatrix}
\]
have linear rank three but require four shared mixed labels.  In any
nonnegative rank-one summand, at most one of the four positive diagonal
entries can be covered: covering two would make both corresponding
cross entries positive, whereas at least one is zero.  Four summands
are therefore necessary, and the displayed rows supply a four-label
realization.  This also explains why a separate scalar module bound
for each preparation cannot replace joint positive synthesis.

\subsection{Robust joint certificates and sharp resource separations}
For weights $\alpha_i\ge0$ with $\sum_i\alpha_i=1$, let
\[
 \mathcal A_{\rho,\alpha}(\mathcal B)
   =[\sqrt{\alpha_1}A_\rho(b_1)\ \cdots\
                       \sqrt{\alpha_{|I|}}A_\rho(b_{|I|})].
\]

\begin{theorem}[Joint robust multiplicity certificate]
\label{thm:r77-joint-robust}
If one mixed shared controller approximates every $b_i(\gamma(w))$
to total-variation error at most $\delta$ for every word $w$, there
is a $G$-set $X$ of size at most its memory such that, for every
irreducible $\rho$ and every such weight vector,
\begin{equation}\label{eq:r77-weighted-fourier}
 \sigma_{m_\rho(X)+1}
       (\mathcal A_{\rho,\alpha}(\mathcal B))\le\sqrt2\delta.
\end{equation}
Hence the integer cover \eqref{eq:r77-joint-module-cover}, with each
$r_\rho$ replaced by the number of singular values strictly above
$\sqrt2\delta$ in any chosen weighted block, is a memory lower bound.
Weights may be chosen separately for different irreducibles.
The same lower bound applies to pure shared controllers.
\end{theorem}
\begin{proof}
Apply Theorem~\ref{thm:r77-shared-normal-form} to the closed convex
total-variation balls.  It supplies a common permutation action and
tables $a_i(g)=\lambda_iP_gD$ with the same error bound.  The calculation
in Theorem~\ref{thm:r76-robust-multiplicity} gives
$\|A_\rho(b_i)-A_\rho(a_i)\|_{\rm op}\le\sqrt2\delta$.
For block rows $z=(z_i)_i$, Cauchy--Schwarz therefore gives
\[
 \left\|\sum_i\sqrt{\alpha_i}
       \{A_\rho(b_i)-A_\rho(a_i)\}z_i\right\|
 \le\sqrt2\delta\sum_i\sqrt{\alpha_i}\|z_i\|
 \le\sqrt2\delta\left(\sum_i\|z_i\|^2\right)^{1/2}.
\]
The weighted joint error has operator norm at most $\sqrt2\delta$.
The approximating joint block has rank at most $m_\rho(X)$ by
Lemma~\ref{lem:r77-joint-multiplicity}; zero weights can only reduce
that rank.  The singular-value nullspace argument used in
Theorem~\ref{thm:r76-robust-multiplicity} proves
\eqref{eq:r77-weighted-fourier}.  Orbit decomposition gives the stated
cover.  The pure normal form has a core no larger than its total memory.
\end{proof}

\begin{proposition}[Unbounded joint multiplicity for a fixed group]
\label{prop:r77-shared-involution}
For every integer $q\ge1$ there is a $q$-preparation experiment on
$G=C_2$ with $q+1$ terminal outcomes such that each preparation alone
has pure and mixed memory two, while their common-device minima are
\[
 \kappa_{C_2}(\mathcal B)=m_{C_2}(\mathcal B)=2q.
\]
Its joint sign multiplicity is $q$, although the sign representation
has degree one.  All target response probabilities may be chosen rational and strictly
positive.  Moreover, in the explicit family below every common
approximation with uniform total-variation error
$\delta<\epsilon/\sqrt{2q}$ still requires at least $2q$ labels.
\end{proposition}
\begin{proof}
Take $C=\{0,1,\ldots,q\}$, the uniform law
$\pi=(q+1)^{-1}\mathbf1$, and $v_i=e_i-e_0$ for $1\le i\le q$.
For any rational $0<\epsilon<1/(q+1)$ put
\[
 b_i(e)=\pi+\epsilon v_i,\qquad
 b_i(s)=\pi-\epsilon v_i,
\]
where $s$ is the nonidentity element of $C_2$.  Each individual response
is nonconstant and is realized by a two-cycle with these two decoder
rows, so its exact memory is two.  For a shared realization the joint
sign block is the $(q+1)\times q$ matrix
$\epsilon[v_1\ \cdots\ v_q]$, of rank $q$.  Every two-cycle in a
$C_2$-set contributes one sign copy; fixed points contribute none.
At least $q$ two-cycles, or $2q$ labels, are necessary.  The disjoint
union of the $q$ displayed two-cycles attains this bound, with pure
preparation $i$ at the positive point of its own cycle.

For equal weights $\alpha_i=1/q$, the squared nonzero singular values
of the joint sign block are the eigenvalues of
$(\epsilon^2/q)(I+\mathbf1\mathbf1^{\mathsf T})$.
Its least singular value is at least $\epsilon/\sqrt q$ (also for
$q=1$).  When $\sqrt2\delta<\epsilon/\sqrt q$, the joint robust
certificate still forces $q$ sign copies.  This lower bound allows
arbitrary stochastic transitions and arbitrary stochastic terminal
decoders satisfying the all-word approximation requirement.
\end{proof}

\begin{proposition}[An unbounded pure initialization cost]
\label{prop:r77-many-initializers}
For every integer $r\ge1$ there is an $r$-preparation binary experiment
on $C_2$ with
\[
 \kappa_{C_2}(\mathcal B)=2,
 \qquad m_{C_2}(\mathcal B)=r+1,
\]
although every preparation separately has pure memory two.
\end{proposition}
\begin{proof}
Choose $0<t_1<\cdots<t_r<1/2$ and set
\[
 b_i(e)=(1/2+t_i,1/2-t_i),\qquad
 b_i(s)=(1/2-t_i,1/2+t_i).
\]
Use one two-cycle with decoder first coordinates $1/2+t_r$ and
$1/2-t_r$.  Appropriate mixtures of its two states realize every
preparation, proving the mixed upper two; nonconstancy proves its lower.
The last preparation uses the positive core point, and the other $r-1$
preparations use distinct transient initialization labels.  This gives
$r+1$ pure labels by Theorem~\ref{thm:r77-shared-normal-form}.

For the lower bound, the $r$ distinct empty-word laws require $r$
distinct pure initial labels.  If those were all the labels, every
decoder row would have first coordinate strictly greater than $1/2$.
Every stochastic transition from any of them would then still give a
first-coordinate probability strictly greater than $1/2$.  This
contradicts the required response after $s$, whose first coordinate is
strictly less than $1/2$.  At least one further label is necessary.
\end{proof}

Both families are ordinary finite causal experiments, and the constructions
also admit diagonal quantum realizations: place the decoder laws on the
diagonal, use permutation unitaries for commands, and use the stated
diagonal preparations.  No quantum advantage is asserted.  Their purpose
is to separate simultaneous preparation, noncyclic multiplicity, and
counted pure initialization from the already solved single-preparation
resource.  For arbitrary different mean laws, arbitrary quantum preparation
sets, or nonreversible commands, the displayed common-mean equality is
not asserted.
